\chapter{Panduan Praktikan}

\section{Ketentuan Umum}

Seluruh pengumpulan bersifat individual. Diskusi konsep selama praktikum
diperbolehkan, tetapi DDL, query, skrip ETL, hasil analisis, dan laporan harus
merupakan pekerjaan sendiri. Kode yang diadaptasi dari sumber lain wajib diberi
sitasi.

Sepuluh modul ini bukan sepuluh latihan yang terpisah. Keluaran satu modul
menjadi masukan modul berikutnya, sehingga pada akhir semester Anda memiliki
satu gudang data yang utuh. Tertinggal satu modul berarti tertinggal pada
seluruh modul sesudahnya --- karena itu setiap modul menyediakan
\berkas{snapshot.sql} berisi keadaan akhir modul sebelumnya bagi yang perlu
mengejar.

\section{Sebelum Praktikum}

Persiapan dilakukan sebelum sesi agar 120 menit tatap muka dapat difokuskan
pada perancangan dan implementasi. Praktikan wajib:

\begin{enumerate}[leftmargin=1.4em]
  \item membaca bagian Konsep Dasar modul yang akan dikerjakan;
  \item menjalankan uji lingkungan pada panduan persiapan;
  \item memastikan dataset berukuran sesuai kebutuhan modul telah tersedia;
  \item mengerjakan butir Pre-lab pada modul; dan
  \item memulihkan \berkas{snapshot.sql} bila luaran modul sebelumnya belum
        lengkap.
\end{enumerate}

Praktikan yang belum menyelesaikan persiapan tetap mengikuti jadwal kelas dan
tidak memperoleh tambahan waktu pada akhir sesi.

\section{Alur Praktikan Selama 120 Menit}

\begin{center}
\small
\begin{tabular}{@{}p{0.62\textwidth}r@{}}
\toprule
\textbf{Kegiatan praktikan} & \textbf{Waktu} \\
\midrule
Mengikuti briefing dan mencatat contoh yang dikerjakan asisten & 15 menit \\
Mengerjakan prosedur terbimbing Bagian A sampai D & 80 menit \\
Menunjukkan luaran saat checkpoint dan memperbaiki butir yang gagal & 15 menit \\
Mencatat tugas luar laboratorium beserta tenggatnya & 10 menit \\
\midrule
\textbf{Total} & \textbf{120 menit} \\
\bottomrule
\end{tabular}
\end{center}

Pada menit ke-95 Anda harus sudah memiliki artefak yang diminta modul agar
skrip pemeriksa dapat dijalankan. Pada menit ke-110, pastikan seluruh skrip,
keluaran, dan catatan telah tersimpan di repositori agar pekerjaan dapat
dilanjutkan secara reproduktif di luar sesi.

\section{Format Pengumpulan}

Seluruh luaran dikumpulkan sebagai repositori Git dengan struktur berikut.

\begin{lstlisting}[style=shellgaya]
praktikum-pgd-<nim>/
  docker-compose.yml
  modul-01-setup/
  modul-02-star-schema/
  ...
  modul-09-olap-dashboard/
  modul-10-kualitas-data/
  proyek-akhir/
  README.md
\end{lstlisting}

Berlaku ketentuan berikut:

\begin{itemize}
  \item penomoran modul memakai dua digit, dari \berkas{modul-01} sampai
        \berkas{modul-10};
  \item nama berkas memakai \perintah{snake\_case} atau tanda hubung, tanpa
        spasi;
  \item laporan mengikuti \berkas{templates/template-laporan.md};
  \item berkas berisi kredensial diberi nama \berkas{.env} dan tidak boleh
        masuk version control;
  \item riwayat commit harus memperlihatkan pekerjaan bertahap, bukan satu
        commit berisi seluruh berkas menjelang tenggat.
\end{itemize}

\section{Peta Modul, Kuliah Penopang, dan CPMK}

Setiap modul menopang capaian pembelajaran mata kuliah tertentu melalui
pertemuan kuliah yang mendahuluinya. Kode CPMK pada tabel berikut juga
tercantum pada tabel Identitas Modul di awal setiap bab.

\begin{center}
\small
\begin{tabular}{@{}cl>{\raggedright\arraybackslash}p{0.34\textwidth}>{\raggedright\arraybackslash}p{0.21\textwidth}@{}}
\toprule
\textbf{Modul} & \textbf{Minggu} & \textbf{Kuliah penopang} & \textbf{CPMK} \\
\midrule
1  & 3  & P2 Desain Database dan Konseptual Gudang Data & CPMK-072 \\
2  & 4  & P3 Desain Logikal Bagian 1 & CPMK-072 \\
3  & 5  & P4 Desain Logikal Bagian 2 & CPMK-092 \\
4  & 6  & P5 Desain Logikal Bagian 3 & CPMK-092 \\
5  & 7  & P6 Desain Fisikal Bagian 1 & CPMK-092 \\
\multicolumn{4}{@{}l}{\textit{Minggu 8 --- Ujian Tengah Semester}} \\
6  & 9  & P7 Desain Fisikal Bagian 2 & CPMK-092 \\
7  & 10 & P9 Integrasi Multi-Sumber dan Rekonsiliasi Skema & CPMK-092 \\
8  & 11 & P10 Proses ETL/ELT dan P11 Metadata, Lineage, dan Auditabilitas
        & CPMK-092 dan CPMK-102 \\
9  & 12 & P12 Pengolahan Data dengan OLAP & CPMK-102 \\
10 & 13 & P13 Manajemen Kualitas Data & CPMK-102 \\
\multicolumn{4}{@{}l}{\textit{Minggu 14--15 --- Proyek akhir: P14 dan P15,
  CPMK-102}} \\
\bottomrule
\end{tabular}
\end{center}

Pertemuan 14 dan 15 tidak memperoleh modul tersendiri. Keduanya dinilai secara
integratif lewat proyek akhir: P14 melalui kelengkapan paket dokumentasi dan
bukti kerja bertahap pada riwayat Git, dan P15 melalui runbook operasional
beserta bukti satu siklus muat berulang yang idempoten.

\section{Prinsip Evaluasi}

Nilai tidak ditentukan oleh skrip yang berjalan saja. Ketepatan grain, alasan
pemilihan tipe data, urutan kolom indeks, jenis fact table, dan pilihan agregat
memiliki bobot penting. \textbf{Jawaban benar tanpa alasan memperoleh separuh
nilai.}

\begin{center}
\small
\begin{tabular}{@{}p{0.34\textwidth}rp{0.42\textwidth}@{}}
\toprule
\textbf{Komponen} & \textbf{Bobot} & \textbf{Dasar penilaian} \\
\midrule
Checkpoint lab, 10 modul & 25\% & Skrip pemeriksa otomatis, lulus atau gagal
  per butir \\
Tugas luar lab, 10 tugas & 25\% & Ketepatan dan alasan \\
Proyek akhir & 40\% & Kelengkapan artefak dua proses bisnis \\
Keaktifan dan disiplin repositori & 10\% & Keteraturan commit, kejelasan
  \berkas{README} \\
\bottomrule
\end{tabular}
\end{center}

Rubrik rinci per modul tersedia pada \berkas{rubrik/rubrik-modul.md}.

\section{Larangan}

\begin{itemize}
  \item menyalin skrip atau laporan praktikan lain;
  \item mengubah data pada basis data sumber untuk membuat uji kualitas lulus;
  \item melaporkan angka hasil query tanpa menyertakan query yang
        menghasilkannya;
  \item memakai angka perkiraan jumlah baris ketika modul meminta
        \perintah{COUNT(*)} sebenarnya;
  \item menuliskan kredensial di dalam skrip yang dikumpulkan; dan
  \item mengumpulkan skrip yang hanya berjalan pada path atau akun pribadi.
\end{itemize}
